\paragraph{線形代数の練習問題}行列$\bm{A}=\begin{pmatrix}1 & 2 & -1\\3 & 1 & 0\end{pmatrix}$のとき，次の計算をしなさい。なお，$\bm{I}_n$とは$n \times n$の単位行列，$\bm{O}$とはすべての要素が0の適当なサイズの正方行列であることを表します。
\begin{enumerate}
    \item $\bm{A}'\bm{A}$
    \item $\bm{AA}'$
    \item $\bm{A}\bm{I}_3$
    \item $\bm{A}$
\end{enumerate}
\paragraph{線形代数の練習問題その2}行列$\bm{A}=\begin{pmatrix}1 & 3 \\2 & 4\end{pmatrix},\bm{B}=\begin{pmatrix}3&1&5\\2&4&6\end{pmatrix},\bm{C}=\begin{pmatrix}4&1\\2&3\end{pmatrix}$，列ベクトル$\bm{x}=\begin{pmatrix}1\\2\\3\end{pmatrix}$,行ベクトル$\bm{y}=\begin{pmatrix}2&8\end{pmatrix}$とするとき，次の計算をしなさい。なお，計算が定義されていないものについては「計算できない」と回答しなさい。
